\section{General Discussion and Limitations}
\label{sec:discussion}

\subsection{Implications for structural assessment}

The four research questions jointly define how structural differences are
represented, compensated, scored, and localized. Low-resource parsing must
be evaluated on the generalization axis relevant to the application,
because the best in-corpus parser did not transfer best to unseen
characters. These rankings describe the tested training protocols,
including their different initialization schemes. Overlapping labels retain
crossing evidence that a mutually exclusive pixel partition cannot express
(Fig.~\ref{fig:segmentation-qualitative}).

Registration shapes the measurement because the transform family
determines which discrepancies remain visible. The bounded search reduces
global nuisance penalties while retaining penalties for structural edits.
The increase in structural penalties can partly reflect centroid-based
refitting after a local edit; it does not independently establish better
measurement fidelity. The tested wider search also uses coarser sampling,
so its comparison does not isolate transform range. Residual sensitivity
prevents a claim of complete invariance under either search.

Scalar agreement and localized diagnostics require different evidence.
ASDS depends on the union silhouette, so permuting direction labels leaves
it unchanged. Such a change can alter directional comparisons and
localized attributions while leaving the separate endpoint mask unchanged.
The direct-ink control supports thresholded ink as an input for scalar
silhouette comparison in this cohort. Within this framework, the parser
supplies direction-specific and endpoint evidence unavailable from the
silhouette alone.

The frozen ASDS specification now requires evaluation without modification
on new writers, characters, and raters. Diagnostic localization also
requires region-overlap or pixel-overlap metrics that credit multi-cell
truth.

\subsection{Limitations}

The segmentation corpus contains 40 labelled character identities from one
recovered acquisition source. Character-disjoint evaluation does not
establish writer-disjoint generalization (Sect.~\ref{sec:data-corpus}), and
collection-specific regularities may survive semantic quality control.

Controlled perturbations operate on cached reference masks predicted by
the v1 checkpoint, which predates semantic quality control and the formal
comparison. Systematic parsing errors in those masks can propagate into
RQ2 and RQ4 statistics. Perturbed masks are compared directly without
rerunning the parser, so this benchmark does not measure end-to-end
performance on perturbed images. Scale-up and compound summaries are
conditional on clipping-safe references because their invalid fractions are
high. The external library supports only seven same-character cross-style
pairs, limiting conclusions about discrimination across styles.

The human cohort contains 150 pairs rated by three domain-qualified raters,
two of whom are authors. Single-rater reliability is poor, while reliability
of the three-rater mean is moderate. The primary score comparison uses
annotated masks, so it excludes deployment parser error; the direct-ink
control bypasses parsing. ASDS feature design followed inspection of the
ratings, making its association retrospective and its character-grouped
validation internal. The direct-ink control reuses this cohort and provides
no independent confirmation.

The contribution of registration to ASDS was not isolated. Its polar term
is centroid- and scale-normalized, so part of the observed association may
persist without alignment. Agreement at multi-label crossing pixels is not
measured separately, and diagnostic ground truth comes from engineered
interventions rather than teacher judgments. The supported modality remains
clean direct digital writing; robustness to camera capture has not been
evaluated (Sect.~\ref{sec:introduction}).
